% ===================================================================
%  नकसा नं. ४ — भरल उमिर आ बुढ़ापा।
%
%  Traduction, faite en 2026, en bhojpouri d'aujourd'hui, sur l'ido
%  de texto/io. Regles de la colonne : voir 10-nakasa-01.tex. DEUX
%  scenes, deux numerotations qui repartent de 1 (t04-c1-*,
%  t04-c2-*). 43 blocs, 93 renvois.
%
%  LE BHOJPOURI NE PEUT PAS DIRE « BEAU- » : IL DOIT SAVOIR DE QUEL
%  COTE, ET LE TRADUCTEUR EST DONC OBLIGE DE LIRE LA PLANCHE. Le
%  francais et l'ido ont un prefixe qui sert a tout — bopatrulo,
%  bomatro, bofilio, bofiliino, bofrato, bofratino — et une meme
%  forme y couvre plusieurs personnes. Le bhojpouri n'a pas de
%  prefixe : il a un mot par place, et chaque mot dit de quel cote de
%  l'alliance on parle. Il faut donc trancher a chaque fois :
%
%    bopatrulo   (20)  ससुर    le pere de l'epouse
%    bomatro     (21)  सास     la mere de l'epoux
%    bofiliino         पतोह    la femme du fils
%    bofilio           दामाद   le mari de la fille
%    bofratino   (26)  ननद     la soeur du mari
%    bofrato     (30)  साढू    le mari de la soeur de l'epouse
%
%  LE DERNIER EST LE PLUS INSTRUCTIF. L'ido dit « la bofrato di
%  Ioannes », ce qui en francais couvre quatre hommes differents ;
%  mais le meme alinea precise qu'il est assis a cote de la
%  « seniora fratino di Martha », donc qu'il a epouse la soeur ainee
%  de l'epouse. En bhojpouri cela ne peut se dire que d'une facon :
%  साढू. La traduction porte ainsi un renseignement que l'original
%  n'a pas, et il a fallu aller le chercher dans le texte pour ne pas
%  se tromper. UNE LANGUE PLUS PRECISE QUE SA SOURCE OBLIGE A MIEUX
%  LIRE LA SOURCE.
%
%  LE CALENDRIER (43) N'EST PAS UN « पतरा », ET C'EST LE TROISIEME
%  REFUS DE LA COLONNE. Le पतरा est l'almanach luni-solaire de la
%  plaine : il donne les tithis, les nakshatras, les jours fastes. Ce
%  que le grand-pere regarde au mur donne le 1er decembre. On ecrit
%  donc « कैलेंडर ». Le tableau 3 avait deja refuse चइत pour les mois
%  et राम-राम pour le salut. Trois refus, une seule loi : ON
%  MODERNISE LA LANGUE, JAMAIS LES CHOSES — et c'est toujours le mot
%  le plus juste de la langue qui est ecarte, jamais le plus
%  lointain.
%
%  DIX-HUIT PAIRES A RETOURNER, ET CETTE FOIS LA LISTE EST RESTEE
%  OUVERTE BLOC PAR BLOC. Le tableau 3 en comptait vingt-deux et en
%  avait manque six pour avoir lu parigi.py d'un coup au lieu de le
%  tenir. Ici l'ordre a ete verifie bloc a bloc pendant l'ecriture,
%  pas apres : la perdrix (2) devant le garde (5), le buisson (4)
%  devant le chien (6), la bouteille (52) devant le tire-bouchon
%  (51), le champagne (49) devant les assiettes (47), l'aveugle (20)
%  devant la pharmacie (16), la grand-mere (23) devant son lorgnon
%  (26), le coussin (38) devant le tabouret (39).
% ===================================================================

%%K t04-apar-1 apar
\VUsaut{4.0mm}
\VUtitre{15.0pt}{60}{नकसा नं. ४}
\VUsaut{1.6mm}
\VUtitre{11.4pt}{0}{भरल उमिर आ बुढ़ापा।}
\VUsaut{3.2mm}

%%K t04-tit-1 sub
\VUcentre{11.4pt}{0}{\textit{पहिलका दृस्य।}}
\VUsaut{1.6mm}
\VUtitre{11.4pt}{0}{बियाह के भोज।}
\VUsaut{2.6mm}

%%K t04-tit-2 sub
\VUpk{10.2pt}{40}{पतझड़।}
\VUsaut{3.0mm}

%%K t04-c1-01-1 p
१. --- जवान लोग बड़ हो गइल। ओहमें से एगो, योहानेस, बियाह करत बा। हमनी
के ओह \VUgras{बियाह के भोज} में बानी जा जवन एके मेज के चारो ओर दुनो
पच्छ के परिवार के जुटवले बा।

%%K t04-c1-02-1 p
२. --- हमनी के \VUgras{पतझड़} में बानी जा, \VUgras{सितंबर महीना} में।
\VUgras{अक्टूबर} आ \VUgras{नवंबर} के ठंढा हवा अबहीं देहात से ओकर हरियर
पत्ता ना छीनले बा। साँझ के हवा गुनगुन आ महकत बा ; एहीसे खाना बगइचा
लगे \VUgras{बरामदा} \textsuperscript{(8)} के नीचे होत बा।

%%K t04-c1-03-1 p
३. --- दूर, \VUgras{पहाड़ी} \textsuperscript{(3)} पर, योहानेस के माई-बाबूजी
के घर बा, एगो सुंदर \VUgras{किला} \textsuperscript{(1)} जवन हरियरी में
छिपल बा ; सुंदर अंगूर के बगइचा, जवना से बढ़िया सराब बनेला, ओह छोट
पहाड़ी के ढलान के ढँकले बा। अंगूर के किसान ओकर जतन करेला।

%%K t04-c1-04-1 p
४. --- अंगूर के बगइचा के ऊपर से \VUgras{तीतर} \textsuperscript{(2)} के
झुंड उड़त बा, जे \VUgras{सिकार-रखवार} \textsuperscript{(5)} के लगे आवे
से डेराइल बा। ऊ, बगल में \VUgras{बंदूक} आ कान्ह पर सिकार-थैला लटकवले,
\VUgras{झाड़ी} \textsuperscript{(4)} में, आपन \VUgras{कुकुर}
\textsuperscript{(6)} के साथे, खोजत बाड़ें। ऊ जमीन के रखवारी करेलें आ चोर
सिकारी के सिकार उजाड़े ना देलें।

%%K t04-c1-05-1 p
५. --- \VUgras{बरामदा} के बाहर \VUgras{अंगूर के बेल} चढ़ल बा, जवना पर
से घन \VUgras{अंगूर के गुच्छा} \textsuperscript{(7)} लटकत बा।

%%K t04-c1-06-1 p
६. --- \VUgras{मेज} \textsuperscript{(16)} के सजावे वाला सुंदर पतझड़ के
फूल \VUgras{गुलदाउदी} \textsuperscript{(9)} ह ; दुलहिन के \VUgras{बाबूजी}
\textsuperscript{(10)} ओहमें से एगो आपन कोट के काज में लगवले बाड़ें।

%%K t04-c1-07-1 p
७. --- खाना खतम होत बा। \VUgras{नौकर} \textsuperscript{(11)} मिठाई के
थारी ले आवे लागल बा आ थोड़े देर में ऊ खाना के ऊ सामान उठा ली जेकर आउर
जरूरत नइखे — \VUgras{चम्मच} \textsuperscript{(14)}, \VUgras{काँटा}
\textsuperscript{(12)}, \VUgras{छूरी} \textsuperscript{(13)},
\VUgras{बड़का थारी} \textsuperscript{(15)}।

%%K t04-tit-3 sub
\VUpk{10.2pt}{40}{नात-रिस्तेदार।}
\VUsaut{3.0mm}

%%K t04-c1-08-1 p
८. --- सब लोग खुस लागत बा, आ दुलहा के \VUgras{माई} \textsuperscript{(17)}
जवन मुसकी से आपन बेटा-बेटी के निहारत बाड़ी ऊ ओकर खुसी के परमान ह।

%%K t04-c1-09-1 p
९. --- ई बियाह इलाका के दू गो धनी परिवार के जोड़त बा। योहानेस,
\VUgras{दुलहा} \textsuperscript{(18)}, एगो बड़का जमींदार के \VUgras{बेटा}
ह, आ ऊ एगो होसियार कारखानेदार के \VUgras{दामाद} बनत बा।

%%K t04-c1-10-1 p
१०. --- \VUgras{दुलहा-दुलहिन} जवान बाड़ें आ तबो ओह लोग के \VUgras{सगाई}
बहुते पहिले हो गइल रहे। \VUgras{जवान दुलहिन} \textsuperscript{(19)},
मार्था, आपन \VUgras{उजर पोसाक} में, जवन संतरा के फूल से सजल बा,
मनमोहक लागत बाड़ी।

%%K t04-c1-11-1 p
११. --- ओकर \VUgras{ससुर} \textsuperscript{(20)} एतना भल \VUgras{पतोह}
पाके बहुते खुस बाड़ें, आ योहानेस के \VUgras{सास} \textsuperscript{(21)}
आपन \VUgras{बेटी} के एतना सुंदर देख के मुसकात बाड़ी।

%%K t04-c1-12-1 p
१२. --- मेज के छोर पर बाकी \VUgras{नेवतहरी} \textsuperscript{(22)} लोग
बइठल बा — \VUgras{नात, संगी, चचेरा भाई, चचेरी बहिन}।
\VUgras{भोज-परोसनिहार} \textsuperscript{(23)}, बगल में नैपकिन दबवले, ओह
लोग के फुरती से परोसत बा।

%%K t04-tit-4 sub
\VUpk{10.2pt}{40}{दोस्त लोग।}
\VUsaut{3.0mm}

%%K t04-c1-13-1 p
१३. --- मेज के दाहिने ओर, इहाँ एगो \VUgras{कुँआरी लइकी}
\textsuperscript{(24)} बाड़ी, जवान दुलहिन के \VUgras{सहेली}, आ ओकरा बाद
ओकर भाई, जे ओकर \VUgras{सम्मान-संगी} \textsuperscript{(25)} ह। ऊ आपन
\VUgras{सम्मान-सहेली} \textsuperscript{(26)} से बतियावत बा, जे दुलहा के
बहिन ह आ जे एह बियाह से मार्था के \VUgras{ननद} बनत बाड़ी। ऊ मनमोहक
जवान बाड़ी। ओकरा लगे सुंदर \VUgras{गहना}, \VUgras{मोती के माला} आ सोना
के \VUgras{कंगन} बा।

%%K t04-c1-14-1 p
१४. --- ओह लोग लगे जे साहेब ठाढ़ होके आपन गिलास उठावत बाड़ें ऊ परिवार
के \VUgras{दोस्त} \textsuperscript{(27)} हउवें। ऊ \VUgras{दुलहा के गवाह}
में से एगो हउवें। ऊ \VUgras{जाम उठावत} बाड़ें, जवान जोड़ी के
\VUgras{सेहत} खातिर पियत बाड़ें आ ओह लोग के लमहर खुसी के कामना करत
बाड़ें। ऊ रस्मी पोसाक पहिनले बाड़ें, कोट आ \VUgras{पॉलिस कइल जूता}
\textsuperscript{(28)} के साथे। ऊ \VUgras{आजी} \textsuperscript{(29)} के
साथे आइल बाड़ें, जे \VUgras{बिधवा} बाड़ी।

%%K t04-c1-15-1 p
१५. --- \VUgras{कोट} \textsuperscript{(31)} पहिनले ऊ जवान, पातर करिया
मोंछ वाला, योहानेस के \VUgras{साढू} \textsuperscript{(30)} हउवें। ऊ एगो
नामी इंजीनियर हउवें। ऊ बहुते सजल-धजल \VUgras{जवान मेहरारू}
\textsuperscript{(33)} के बगल में बाड़ें, जे मार्था के \VUgras{बड़की
बहिन} बाड़ी आ ओह नन्हीं लइकी के माई बाड़ी जे आपन चचेरा भाई के हाथ पकड़
के फूल देवे आ ओकरा के \VUgras{नीक साँझ} \VUgras{कहे} आवत बाड़ी। ओह
मेहरारू लगे बहुते सुंदर \VUgras{कान के बाली} आ सुंदर \VUgras{केस-सिंगार}
बा।

%%K t04-c1-16-1 p
१६. --- छोटका चचेरा भाई, आपन \VUgras{बुसकट} \textsuperscript{(36)} आ आपन
फुलल \VUgras{जाँघिया} \textsuperscript{(37)} के चलते एतना अजीब लागत बा,
बहुते बेचारा बा। ऊ \VUgras{अनाथ} \textsuperscript{(35)} ह। बहुते छोट में
ऊ आपन बाबूजी आ आपन माई के खो देलस। ओकर चाचा ओकर \VUgras{अभिभावक}
\textsuperscript{(38)} हउवें, आ ऊ एह बेचारा बच्चा के पाले खातिर कुँआर
रह गइलें। ऊ भल मनई हउवें। ऊ थोड़ा गंजा बाड़ें आ बहुते कम नजर वाला ;
ऊ \VUgras{चस्मा} लगावेलें।

%%K t04-tit-5 sub
\VUsaut{2.4mm}
\VUpk{10.2pt}{40}{मिठाई-दौर।}
\VUsaut{3.0mm}

%%K t04-c1-17-1 p
१७. --- भोज करे वाला लोग के पाछे, \VUgras{मेजपोस} से ढँकल एगो मेज पर,
\VUgras{मिठाई-दौर} \textsuperscript{(39)} तइयार बा। बड़का कटोरा में फर
बा : \VUgras{आड़ू} \textsuperscript{(40)}, \VUgras{नासपाती}
\textsuperscript{(41)}, \VUgras{सेब} \textsuperscript{(42)},
\VUgras{खुबानी} \textsuperscript{(43)} आ \VUgras{अंगूर} \textsuperscript{(44)}।
ओकरा लगे \VUgras{अनन्नास} \textsuperscript{(45)}, सुंदर \VUgras{केक}
\textsuperscript{(46)}, \VUgras{सराब के बोतल} \textsuperscript{(48)} आ
\VUgras{सैंपेन} \textsuperscript{(49)}, आ साफ \VUgras{थारी}
\textsuperscript{(47)} के थाक भी बा।

%%K t04-c1-18-1 p
१८. --- \VUgras{सराब-परोसनिहार} \textsuperscript{(50)} पुरनका सराब के
\VUgras{बोतल} \textsuperscript{(52)} \VUgras{पेंच} \textsuperscript{(51)}
से खोलत बा। \VUgras{कॉर्क के डाट} \textsuperscript{(53)} बहुते मुसकिल से
निकलत लागत बा। ओह परोसनिहार के बगल में \VUgras{नैपकिन}
\textsuperscript{(54)} दबल बा।

%%K t04-c1-19-1 p
१९. --- खाना के बाद साहेब लोग सिगरेट पिए के कमरा में जइहें, आ मेहरारू
लोग नाच खातिर तइयार होखे जइहें।

%%K t04-tit-6 sub
\VUsaut{5.0mm}
\VUcentre{11.4pt}{0}{\textit{दुसरका दृस्य।}}
\VUsaut{1.6mm}
\VUpk{11.4pt}{40}{बाबा के जनम-दिन।}
\VUsaut{2.6mm}

%%K t04-tit-7 sub
\VUpk{10.2pt}{40}{भेंट।}
\VUsaut{3.0mm}

%%K t04-c2-01-1 p
१. --- दस बरिस बीत गइल, योहानेस के माई-बाबूजी बूढ़ हो गइलें आ आपन तीन
गो नाती-नतिनी से बहुते मया करेलें — दू गो \VUgras{नाती}
\textsuperscript{(2)} आ एगो \VUgras{नतिनी} \textsuperscript{(3)}।
काहेकि आज \VUgras{बाबा के जनम-दिन} बा, बेटा-बेटी लोग ओह लोग के
बधाई देवे आइल बा।

%%K t04-c2-02-1 p
२. --- छोटका पेत्रुस अबहीं बहुते छोट बा, ऊ खाली तीन बरिस के बा। ऊ
\VUgras{बिलार} \textsuperscript{(1)} के लगे आवत देखत बा। ऊ ओकर सुंदर
रेसमी \VUgras{रोआँ} आ ओकर लमहर \VUgras{पूँछ} सहलावल चाहत बा ; ऊ ई ना
सोचत कि ओह सुंदर \VUgras{पंजा} के नीचे तेज नाखून छिपल बा जवन सायद ओकरा
के नोच ली।

%%K t04-c2-03-1 p
३. --- ओकर बहिन गाब्रिएला, जे ओकर हाथ पकड़ले बाड़ी, बहुते जादे
गंभीर बाड़ी, आ मार्केल्लुस आपन बड़का \VUgras{लबादा} \textsuperscript{(4)}
आ चमड़ा के \VUgras{पेटी} \textsuperscript{(5)} में थोड़ा बड़ लागत बा,
हालाँकि ओकर छोट जाँघिया से ओकर \VUgras{मोजा} \textsuperscript{(6)} लउकत
बा।

%%K t04-c2-04-1 p
४. --- ओह लोग के बाबूजी आपन मोट जाड़ के \VUgras{ओवरकोट}
\textsuperscript{(7)} पहिनले बाड़ें \VUgras{फर के कॉलर} \textsuperscript{(8)}
के साथे, आ आपन \VUgras{जेब} \textsuperscript{(9)} में ओकरा लगे रूमाल बा।
ओकर घरवाली लगे \VUgras{पंख} \textsuperscript{(10)} से सजल फेल्ट के टोपी,
आपन \VUgras{जैकेट} \textsuperscript{(11)}, आपन \VUgras{बोआ-फर}
\textsuperscript{(12)} आ आपन \VUgras{मफ} \textsuperscript{(13)} बा।

%%K t04-c2-05-1 p
५. --- बहुते ठंढा बा, काहेकि जाड़ के राज बा। \VUgras{बरफ}
\textsuperscript{(19)} पूरा दिन गिरत रहल आ घरन के छत पूरा उजर बा।
\VUgras{रात} के साथे ठंढ आउर तेज हो जात बा। आसमान में \VUgras{चान}
\textsuperscript{(17)} आ \VUgras{तारा} \textsuperscript{(18)} चमकत
बाड़ें आ \VUgras{अन्हार} के चीरत बाड़ें।

%%K t04-tit-8 sub
\VUsaut{4.2mm}
\VUpk{10.2pt}{40}{बेमारी।}
\VUsaut{3.0mm}

%%K t04-c2-06-1 p
६. --- बेचारा \VUgras{अन्हरा} \textsuperscript{(20)}, जे
\VUgras{दवाखाना} \textsuperscript{(16)} के सामने के चउक आपन \VUgras{पूडल
कुकुर} \textsuperscript{(21)} के सहारे पार करत बा, बहुते दुखी बा।
इउस्तीनिया, \VUgras{नौकरानी} \textsuperscript{(14)}, रसोई से बहुते गरम
\VUgras{काढ़ा} के \VUgras{कटोरा} \textsuperscript{(15)} ले आवत बाड़ी, जवना
में खूब चीनी मिलल बा।

%%K t04-c2-07-1 p
७. --- \VUgras{आजी} \textsuperscript{(23)}, जे मेज पर आपन \VUgras{ऐनक}
\textsuperscript{(26)} खोजल चाहत बाड़ी, ताकि ऊ \VUgras{नुसखा}
\textsuperscript{(27)} पढ़ सकीं जवन \VUgras{डाक्टर} \textsuperscript{(25)}
अभिए लिखले बाड़ें, आपन \VUgras{लाठी} \textsuperscript{(24)} पर टेक के आपन
आरामकुरसी से उठत बाड़ी।

%%K t04-c2-08-1 p
८. --- अकसर ओकरा \VUgras{जी मिचलाइल, चक्कर, सरदी, दाँत के दरद} होला,
ओकर गोड़ कमजोर बा आ ओकर आँख आउर जादे नीक नइखे। तबो ऊ आपन पुरनका
\VUgras{डाक्टर} के हँसी उड़ावेली, जे हरमेसा ओकरा के \VUgras{दवाई}
\textsuperscript{(28)} लिखल चाहेलें। आज ऊ बहुते खुस बाड़ी काहेकि आपन
जवान परिवार ओकरा चारो ओर बा।

%%K t04-c2-09-1 p
९. --- नीक से बंद कमरा में आराम बा। संगमरमर के \VUgras{अँगीठी}
\textsuperscript{(29)} में \VUgras{सुंदर आग} \textsuperscript{(30)} जरत
बा, आ ओह \VUgras{लालटेन} \textsuperscript{(31)} के मीठ \VUgras{अँजोर} में,
जवन अभिए जरावल गइल बा, मन खुस हो जात बा।

%%K t04-tit-9 sub
\VUsaut{4.2mm}
\VUpk{10.2pt}{40}{बाबा।}
\VUsaut{3.0mm}

%%K t04-c2-10-1 p
१०. --- खुद \VUgras{बाबा} \textsuperscript{(33)} भुला गइल बाड़ें कि ऊ
बेमार रहलें, कि ओकर \VUgras{गठिया} ओकरा के आपन \VUgras{आरामकुरसी}
\textsuperscript{(34)} से बान्हले बा, आ कि कल्हुए ओकरा \VUgras{बोखार आ
बेहोसी} रहे। ओकरा ई भी ना इयाद बा कि पिछला जाड़ में ऊ \VUgras{निमोनिया}
से पीड़ित रहलें आ खरखर आ थकावे वाला \VUgras{खाँसी} ओकरा के बहुते सतवले
रहे।

%%K t04-c2-10-2 p
आ तबो ऊ बहुते मुसकिल से नीक भइलें। अबहीं ऊ थोड़ा बेहतर बाड़ें। बाकिर ओकर
जवन गोड़ ऊ \VUgras{कुसन} \textsuperscript{(38)} पर टेकले बाड़ें, जवन छोट
\VUgras{चौकी} \textsuperscript{(39)} पर धइल बा, ऊ भी ओकरा के दुखावेला।

%%K t04-c2-11-1 p
११. --- जब ऊ आपन गरम \VUgras{ड्रेसिंग गाउन} \textsuperscript{(35)} में
सलीका से लपेटल रहेलें, जवना में सीप के मोटका \VUgras{बटन}
\textsuperscript{(36)} लागल बा, आ जब ओकरा लगे अखबार पढ़े खातिर छोटकी
\VUgras{अनाथ लइकी} \textsuperscript{(40)} रहेली, जे उजर \VUgras{टोपी},
नीला \VUgras{फ्राक} \textsuperscript{(42)} आ \VUgras{पेलरिन}
\textsuperscript{(41)} पहिनले बाड़ी, त ऊ कवनो सिकायत ना करेलें।

%%K t04-c2-12-1 p
१२. --- हर साल, जब \VUgras{कैलेंडर} \textsuperscript{(43)} फेरु पहिला
\VUgras{दिसंबर} के \VUgras{तारीख} देखावेला, त ऊ बेसबरी से आपन
\VUgras{नाती-पोता} के अगोरेलें, काहेकि ऊ जानेलें कि ऊ लोग ओह जरूरी
\VUgras{तारीख} के ना भुलाई।

%%K t04-c2-13-1 p
१३. --- ऊ भावुक होके ओह बहुते पुरनका बेरा के इयाद करेलें जब ऊ खुद
सम्राट के नामी \VUgras{मार्शल} के बधाई देवे गइल रहलें, जेकर
\VUgras{तसवीर} \textsuperscript{(44)} भीत पर टँगल बा आ जे ओकर बेटा-बेटी
के \VUgras{परबाबा} हउवें।
